\documentclass[11pt]{article}
\usepackage[a4paper,margin=1in]{geometry}
\usepackage{amsmath,amssymb,amsthm}
\usepackage{booktabs}
\usepackage{graphicx}
\usepackage{xcolor}
\usepackage{hyperref}
\hypersetup{colorlinks=true,linkcolor=blue,urlcolor=blue,citecolor=blue}
\newcommand{\rt}[1]{\texttt{#1}}
\newcommand{\fn}[1]{\texttt{#1}}
\providecommand{\ket}[1]{|#1\rangle}
\newcommand{\wcc}{\mathrm{wcc}}
\newcommand{\rec}{\mathrm{rec}}
\newcommand{\ifgraphic}[2]{\IfFileExists{#1}{\includegraphics[width=#2]{#1}}%
  {\fbox{\parbox{#2}{\centering\small figure \texttt{\detokenize{#1}} pending}}}}
\newcommand{\iftable}[1]{\IfFileExists{#1}{\input{#1}}%
  {\fbox{\parbox{0.9\textwidth}{\centering\small table
  \texttt{\detokenize{#1}} pending}}}}

\theoremstyle{plain}
\newtheorem{theorem}{Theorem}
\newtheorem{proposition}{Proposition}
\newtheorem{corollary}{Corollary}

\title{\textbf{Report R11 --- One rule space, two gates}\\
\large What changes, and what cannot change, when the Hadamard is replaced by
$X$}
\author{QCA Sector Analysis project (Tier 1)}
\date{2026-07-31 \quad\small(R9 \texttt{1e.1} vs R10 \texttt{1x.1}, \rt{obc0},
common window $N\le16$)}

\begin{document}
\maketitle

\begin{abstract}
R1--R9 analyse $256$ brick-wall rules whose active gate is the Hadamard; R10
analyses the same $256$ symbol tables with the gate changed to $X$. This report
compares them. The two families turn out not to be independent: because applying
the Hadamard to a site produces \emph{both} values of the target bit, the flip is
always one of the outcomes, so the $X$ transition graph is a subgraph of the
Hadamard one and the $X$ sectors \textbf{refine} the Hadamard sectors --- verified
exhaustively for all $256$ rules and over $2816$ swept units, with zero
exceptions. Three consequences organise everything else. The gate can only ever
\emph{split} sectors, never merge them, so no rule loses sector structure: of the
$51$ rules with an exponentially growing sector count under $H$, all $51$ keep one
under $X$, and $39$ rules with no structure under $H$ acquire it. The $81$ rules
with no \rt{V} are literally the same circuit under both gates, which makes the
two independent engines a cross-implementation test that they pass on all $916$
shared units. And on the $16$-rule unitary/reversible baseline --- the only fully
like-for-like comparison, since both families satisfy sectors\,$=$\,attractors
there --- the gate change is nearly a reflection: $9$ of the $11$ rules that have
a single giant sector under $H$ end up with no giant sector at all under $X$,
$7$ of them at the opposite extreme of the axis, with rule $204$ the sole fixed point and
rule $51$ going from a single sector of $2^N$ states to $2^{N-1}$ two-cycles.
Quantitatively almost nothing transfers: on the unitary baseline the Pearson
correlation between the two gates' exponents is $+0.17$ in $a$ and $+0.03$ in
$b$. \rt{obc0} only; \rt{pbc} is held for both families.
\end{abstract}

\tableofcontents

\section{What is being compared, and on what window}
\label{sec:method}

The rule encoding is shared exactly. A rule is a tuple
$(r_{00},r_{01},r_{10},r_{11})$ over $\{\rt{I},\rt{V},\rt{D},\rt{E}\}$; the
circuit applies the symbol selected by a site's two neighbours, brick-wall, even
sublattice ascending and then odd, each reading the current state. Both engines
compile the same step list from \fn{core/cycle.py}. The single difference is what
\rt{V} means:

\begin{center}
\begin{tabular}{cll}
\toprule
symbol & Hadamard family (R1--R9) & permutation family (R10) \\
\midrule
\rt{I} & identity & identity \\
\rt{V} & Hadamard $H$ & $X$ \\
\rt{D} & reset to $\ket0$ & reset to $\ket0$ \\
\rt{E} & reset to $\ket1$ & reset to $\ket1$ \\
\bottomrule
\end{tabular}
\end{center}

\paragraph{The window.} R9 fits its bases on $N\le16$ and R10 on $N\le24$. A base
fitted on a longer window is a different number, so every comparison below
rebuilds the X-gate descriptors on \textbf{the same window as the Hadamard ones},
$N\le16$, through the identical \fn{scaling/sectors.py} machinery. Nothing is
compared across windows. Where an $N=24$ number is quoted it is labelled. One
consequence is visible in the tables: rule $108$'s X-gate cycle length is
classified \emph{linear} on $N\le16$ and \emph{exponential} on $N\le24$, which is
a statement about how much data the classifier has, not about the rule.

\paragraph{What is not compared.} The Hadamard family carries objects that need
interference between distinct branches of a basis state --- destructive
cancellation in $\operatorname{succ}$, the R7 visibility theorem's sharp cases.
It does \emph{not} have a monopoly on protected coherence: R10~\S9 finds
decoherence-free subspaces of exponential dimension throughout the permutation
family, and for the $81$ V-free rules they are literally the same subspaces,
since the channel is the same. Conversely the
permutation family's transient structure has no Hadamard counterpart worth
plotting. The comparison is restricted to the two coordinates both families
define: the sector count $n_\wcc\sim a^N$ and the largest sector
$D_{\max}\sim b^N$, plus the growth \emph{class} of each.

\section{The two families are nested}
\label{sec:refine}

The comparison is not between two unrelated data sets. There is an exact
containment, and it is elementary.

\begin{theorem}[refinement]
\label{thm:refine}
For every rule, every $N$ and both boundary conventions,
$\operatorname{succ}_X(x)\in\operatorname{succ}_H(x)$ for every basis state $x$.
Hence the $X$ transition graph is a subgraph of the Hadamard transition graph on
the same vertex set, and the $X$ sectors refine the Hadamard sectors: every
$X$-component lies inside a single $H$-component, and every Hadamard sector is a
disjoint union of $X$ sectors.
\end{theorem}

\begin{proof}
Consider one elementary step. At \rt{I}, \rt{D}, \rt{E} sites the two families do
the same thing. At a \rt{V} site the Hadamard maps $\ket{b}$ to
$(\ket0\pm\ket1)/\sqrt2$, so the measured or Kraus-resolved successor set contains
\emph{both} values of the target bit; the $X$ action, which sets it to $1-b$, is
one of them. The circuit is a composition of such steps, and choosing the flip at
every \rt{V} site along the sweep is one branch of the Hadamard successor set. So
each $X$ edge is an $H$ edge. Union-find over a subgraph produces a refinement.
\end{proof}

\begin{corollary}
\label{cor:mono}
$n_\wcc^X(N)\ge n_\wcc^H(N)$ and $D_{\max}^X(N)\le D_{\max}^H(N)$, pointwise in
$N$; asymptotically $a_X\ge a_H$ and $b_X\le b_H$.
\end{corollary}

\begin{corollary}
\label{cor:nofree}
A rule with no \rt{V} never invokes the gate, so its two circuits are the
\emph{same map}. There are $81$ such rules ($204$ and the $80$ V-free ones), and
for them the refinement is an equality.
\end{corollary}

\noindent All of this is verified rather than assumed:

\begin{center}
\iftable{tab_r11_structure_obc0.tex}
\end{center}

\noindent The last row is worth its own sentence. The Hadamard sectors come from
a set-valued successor relation streamed through union-find; the X-gate sectors
come from a numpy step table and a three-colour walk. The two share no code below
\fn{core/cycle.py::\_compile}. On the $81$ rules where they must agree they agree
exactly --- the full sector-size multiset, at every $N$ where both were run. That
is the strongest cross-implementation check either report has.

\paragraph{Reading the theorem physically.} The Hadamard's extra edge is the
non-flip branch, and it is exactly what \emph{glues} $X$ sectors together. So one
may think of the Hadamard sector structure as the $X$ sector structure after
gluing: superposition is a coarsening operation on the sector lattice. That is
the reverse of the intuition one might start with --- that a superposing gate
would \emph{create} structure --- and it explains the census in the next section
without any further input.

\section{No rule loses its sector structure}
\label{sec:classes}

Corollary~\ref{cor:mono} says the sector count cannot fall. The growth-class
cross-tabulation makes the consequence concrete:

\begin{center}
\iftable{tab_r11_classes_obc0.tex}
\end{center}

\noindent Every entry below the diagonal is zero, as it must be. Reading the rows:

\begin{itemize}
\item \textbf{All $51$ rules with an exponential sector count under the Hadamard
  keep one under $X$.} Hilbert-space fragmentation, measured by the growth class
  of the sector count, is not a property of the Hadamard: it survives replacing
  the superposing gate by a permutation.
\item $39$ of the $173$ rules with a \emph{constant} sector count under $H$ ---
  no structure at all, one or two sectors covering the space --- acquire structure
  under $X$: $28$ exponential, $6$ linear, $5$ irregular. The gate change can only
  add.
\item $10$ of the $30$ rules with a linear sector count under $H$ become
  exponential; the other $20$ stay linear.
\end{itemize}

\noindent The reverse reading is the one that matters for the Hadamard reports:
if a rule has no sector structure under $X$, it has none under $H$ either. The
$134$ rules in the top-left cell are structureless under both gates, and that is a
property of their symbol table rather than of the dynamics.

\begin{figure}[htbp]
\centering
\ifgraphic{../../figures/fig_cmp_all_obc0.pdf}{\textwidth}
\caption{\textbf{C2, all $256$ rules on a common window.} Left: the sector plane,
Hadamard point (circle) joined to $X$ point (square) for each rule that has a
\rt{V}; the $81$ V-free rules are drawn as diamonds, one point each, because the
gate cannot move them (Corollary~\ref{cor:nofree}). Almost every Hadamard point
sits at $(1,2)$ --- one giant sector --- and the gate change fans them out to the
right. Right: the sector-count base under one gate against the other. Nothing lies
below the diagonal, which is Corollary~\ref{cor:mono} drawn.}
\label{fig:c2}
\end{figure}

\section{The unitary / reversible baseline}
\label{sec:baseline}

The $16$ rules built from \rt{I} and \rt{V} alone are the cleanest comparison
available, for a specific reason: on the Hadamard side they are unitary, so
sectors and monitored attractors coincide (R9's check A2), and on the $X$ side
they are reversible, so sectors and cycles coincide (R10's F1). The same object is
being measured in both families, and the comparison cannot be confounded by the
transient structure that dominates everywhere else.

\begin{center}
\small
\iftable{tab_r11_baseline_obc0.tex}
\end{center}

\begin{figure}[htbp]
\centering
\ifgraphic{../../figures/fig_cmp_baseline_obc0.pdf}{\textwidth}
\caption{\textbf{C1, the baseline.} Left: each of the $16$ rules under both
gates, joined by the move; marker area counts rules sharing a cell. Right: the
sector-count base rule by rule. Rule $204$ is the only one the gate leaves where
it was.}
\label{fig:c1}
\end{figure}

\noindent Three things stand out, and none of them is a small effect.

\paragraph{The gate nearly reflects the family.} Eleven of the $16$ sit at
$(a,b)=(1,2)$ under the Hadamard --- a single sector of size $2^N$, no
fragmentation whatever. Under $X$, nine of those eleven end at $b=1$: the largest
sector is no longer a fixed fraction of the space at all. Seven go to $(2,1)$
exactly --- exponentially many tiny orbits, maximal fragmentation --- and $153$ and
$195$ to $(1.83,1)$. Only $57$ and $99$ keep $b=2$, moving left instead of down.
The extreme case is rule $51$
(\rt{VVVV}). Under the Hadamard it is the most ergodic rule in the family, one
sector covering the whole basis. Under $X$ it is $x\mapsto\bar x$, so the space
splits into $2^{N-1}$ two-cycles. Same symbol table, opposite ends of the sector
axis, and Theorem~\ref{thm:refine} says which way round it must be: the Hadamard
picture is the $X$ picture after gluing.

\paragraph{Only $204$ is fixed.} Rule $204$ is \rt{IIII}, the identity, which
contains no \rt{V} and therefore cannot move (Corollary~\ref{cor:nofree}). It is
the only one of the $16$ with that property, and correspondingly the only rule
whose $a$ is identical in the two columns.

\paragraph{The four fragmented rules stay fragmented, with different
constants.} $156$, $198$, $108$, $201$ are the four rules whose sector count grows
exponentially under the Hadamard, and all four still do under $X$ --- but the
constants have no relation:
\[
\begin{array}{lcc}
 & \text{Hadamard} & X \\[2pt]
156,\ 198 & (\varphi,\ 4^{1/5}) & (2,\ 1.298) \\
108,\ 201 & (1.7549,\ \varphi) & (1.828,\ 2),\ (1.847,\ 1.234)
\end{array}
\]
The Hadamard values are exact algebraic numbers with derivations behind them
(R2~\S3's room-packing saddle point for $4^{1/5}$; the golden ratio and the root
of $x^3=2x^2-x+1$ from integer recurrences). The X-gate values are not the same
numbers and are not obviously algebraic. Fragmentation transfers; the
\emph{grammar} producing it does not.

\paragraph{And the correlation is nil.}

\begin{center}
\iftable{tab_r11_correlation_obc0.tex}
\end{center}

\noindent On the $16$ baseline rules the Pearson correlation between the two
gates is $+0.17$ in $a$ and $+0.03$ in $b$, with mean absolute differences of
$0.65$ and $0.63$ on axes of total width $1$. Read the ``all $256$'' row with
care: its apparent $r=+0.64$ is inflated by the $81$ rules that are the same
circuit under both gates and therefore correlate perfectly by construction.
Restricted to the $175$ rules that actually contain a \rt{V}, $r$ falls to
$+0.47$ and $+0.33$. The gate does not preserve exponents.

\section{R9's headline sets under the other gate}
\label{sec:survivors}

R9 singled out three sets of rules. Each is asked here what it does under $X$.

\begin{center}
\small
\iftable{tab_r11_survivors_obc0.tex}
\end{center}

\begin{figure}[htbp]
\centering
\ifgraphic{../../figures/fig_cmp_survivors_obc0.pdf}{\textwidth}
\caption{\textbf{C3, the three sets.} Hadamard (circle) to $X$ (square) in the
sector plane, common window. Labels list the rules sharing a destination cell.}
\label{fig:c3}
\end{figure}

\subsection{The open-system-fragmented rules}

R9's central physical claim is that eight V+reset rules keep an exponentially
growing sector count, and that because a sector is an enclosure --- $\mathrm{span}$
of a weak component is invariant under every Kraus operator --- this is
fragmentation of an \emph{open} quantum system rather than of a unitary circuit
with noise added afterwards.

All eight remain exponential under $X$, and four of them do something stronger:
$140$, $196$, $206$, $220$ have \emph{exactly} the same sector count as under the
Hadamard, $2584$ at $N=16$, because they are V-free and hence the same circuit
(Corollary~\ref{cor:nofree}). Their appearance in R9's list is therefore not a
statement about the Hadamard at all: those four rules are classical, and their
fragmentation is a classical fact that R9 measured with a quantum engine. The
other four --- $156$, $198$, $108$, $201$ --- do contain a \rt{V}, and for them the
$X$ count is strictly larger ($7744$ against $2584$; $16560$ against $10252$;
$12076$ against $5842$), exactly as Corollary~\ref{cor:mono} requires.

This is a sharpening of R9 rather than a contradiction of it. The claim ``these
rules fragment as open quantum systems'' stands; what R11 adds is that for half
of them the Hadamard was never doing any work.

\subsection{The linear-sector-count rules and the pinned frontier}

R9~\S6.4 identified six rules --- $44$, $100$, $110$, $124$, $188$, $230$ --- whose
sector count is strictly linear, $N+1$ of them, carried by a charge that is the
position of the extremal excitation; \S6.5 showed the charge is an \rt{obc0}
artefact, collapsing to two sectors on the ring. Under $X$ the six split three
ways:

\begin{itemize}
\item $188$ (\rt{IIVE}) and $230$ (\rt{IVIE}) keep \emph{exactly} $17$ sectors at
  $N=16$, the same linear count as under the Hadamard, though they contain a
  \rt{V} and are therefore not forced to. For these two the pinned-frontier charge
  is gate-independent.
\item $110$ and $124$ become exponential but only just, $149$ sectors at $N=16$
  against $17$.
\item $44$ and $100$ jump to $595$.
\end{itemize}

\noindent The frontier family of \S6.5 adds $60$ and $102$, and those two go all
the way to $4116$ sectors --- from a linear count to $a=2$, the maximum. So the
pinned frontier is a Hadamard-specific structure for six of the eight rules that
carry it, and survives the gate change for two. Combined with R9's finding that it
does not survive the ring, the charge is doubly fragile: it depends on the
boundary condition \emph{and}, mostly, on the gate. That is a good reason not to
build a physical story on it, and it is the sort of thing the comparison is for.

\section{Reading all of this as Hilbert-space fragmentation}
\label{sec:hsf}

The two coordinates used throughout are exactly the two numbers the HSF
literature uses to classify a fragmented model, so the map is a phase diagram
and it is worth saying so explicitly.

\subsection{The $(a,b)$ plane is the fragmentation phase diagram}

A model is \emph{fragmented} when the Hilbert space --- or a symmetry sector of
it --- splits into exponentially many dynamically disconnected Krylov subspaces.
With $n_\wcc\sim a^N$ subspaces and a largest one of dimension $D_{\max}\sim b^N$:

\begin{center}
\begin{tabular}{lll}
\toprule
$a=1$, $b=2$ & \textbf{unfragmented} & one Krylov sector carries the space \\
$a>1$, $b=2$ & \textbf{weak} & exponentially many, largest still $\sim2^N$ \\
$a>1$, $1<b<2$ & \textbf{strong} & largest sector an exponentially
  vanishing fraction \\
$b=1$ & \textbf{shattered} & every sector sub-exponential \\
\bottomrule
\end{tabular}
\end{center}

\noindent The exclusion curve $ab\ge2$ is then the statement that one cannot have
few sectors \emph{and} small ones, which is why the weak/strong boundary at $b=2$
and the shattered edge at $b=1$ are the two edges of the diagram rather than
independent axes. The corner $a=1$, $b<2$ is forbidden asymptotically, and every
rule the fitter puts there has a sub-exponential sector count that the base cannot
express (R9~\S5); those are labelled ``degenerate'' below and are a fitting
category, not a phase.

\begin{center}
\small
\iftable{tab_r11_phases_obc0.tex}
\end{center}

\noindent Refinement (Theorem~\ref{thm:refine}) forbids the lower-left of this
table too: a gate change to $X$ can only move a rule \emph{towards} more
fragmentation. And it does. Under the Hadamard $50$ of $256$ rules are strongly
fragmented; under $X$, $75$ are, plus $12$ shattered. Not one rule becomes less
fragmented. The two \emph{weak} entries are $57$ and $99$, the only rules anywhere
in either family with exponentially many sectors and a largest sector still of
full exponential size.

\subsection{Rules 156 and 198: a configuration space, or a bundle of orbits}
\label{sec:156}

The clearest single comparison in the report. Rules $156$ (\rt{IIVI}) and $198$
(\rt{IVII}) are strongly fragmented under \emph{both} gates --- and mean something
completely different by it.

\paragraph{Under the Hadamard: a constrained configuration space.} R2~\S3 derives
the structure. Domain walls are frozen; the chain decomposes into ``rooms''
between them; a room of length $\ell$ contributes $\ell+1$ basis states and
consumes $\ell+2$ sites. Counting the wall patterns gives $n_\wcc\sim\varphi^N$
and maximising the room packing gives $D_{\max}\sim(4^{1/5})^N$, the saddle point
of $b(\ell)=(\ell+1)^{1/(\ell+2)}$. So a sector is labelled by a \emph{static}
charge --- the wall pattern --- and inside it the accessible states form an
exponentially large constrained configuration space, on which the Hadamard acts
as a genuine quantum dynamics.

\paragraph{Under $X$: the same sectors, cut into orbits by extra charges.} The
sectors cannot be finer than the Hadamard ones by anything other than refinement,
and indeed the wall pattern is still conserved. But the $X$ dynamics is a
permutation, so each Hadamard sector decomposes further --- and it does so with a
very particular arithmetic. At $N=14$ rule $156$'s $987$ Hadamard sectors split
into $2280$ $X$ orbits, and \textbf{every sector splits into orbits that are all
the same length}:
\[
\dim(\text{$H$ sector}) \;=\; (\text{number of $X$ orbits})\times(\text{orbit length}).
\]

\begin{center}
\small
\iftable{tab_r11_split156_obc0.tex}
\end{center}

\noindent Read the arithmetic. A sector of dimension $16$ appears as $2\times8$,
$4\times4$ or $8\times2$; one of dimension $36$ as $3\times12$ or $6\times6$; one
of dimension $64$ as $16\times4$. The orbit lengths that occur are $1$ through
$15$, then $18$, $20$, $21$, $22$, $24$, $28$, $30$, $35$, $36$, $40$, $42$ and
$60$ --- and $35=\mathrm{lcm}(5,7)$, $42=\mathrm{lcm}(6,7)$,
$60=\mathrm{lcm}(4,15)$. At $N=24$ the longest orbit in the whole rule is $420$,
and the sequence leading there is $105,140,140,210,210,420,420,420$, i.e.
$\mathrm{lcm}(3,5,7)$, $\mathrm{lcm}(4,5,7)$, $\mathrm{lcm}(2,3,5,7)$,
$\mathrm{lcm}(3,4,5,7)$.

That is the signature of \emph{independent periodic subsystems}. The natural
reading, consistent with every number above though not derived here: inside a
fixed wall pattern each room carries its own cyclic degree of freedom of some
period $p_r$; the $X$ dynamics advances all of them together; so the sector
dimension is $\prod_r p_r$, an orbit has length $\mathrm{lcm}_r(p_r)$, and the
number of orbits is $\prod_r p_r/\mathrm{lcm}_r(p_r)$ --- the \emph{relative
phases} between rooms. Those relative phases are exactly the extra conserved
charges that the $X$ gate has and the Hadamard has not: the Hadamard superposes
the room states and destroys the phase label, gluing the orbits back into one
sector.

The property is specific, not generic: it holds for $156$ and $198$ at every $N$
checked, and fails for $108$, $201$, $54$, $150$ and $105$, whose sectors contain
orbits of differing length.

\subsection{Why the Hadamard result is the interesting one}

Three reasons, in increasing order of weight.

\paragraph{The sectors have room in them.} Strong fragmentation is only a
statement worth making if a sector is big enough to have its own thermodynamics.
Rule $156$ under the Hadamard has $\varphi^N$ sectors of size up to
$(4^{1/5})^N$ --- both exponential, product $2.135>2$ --- so each Krylov sector is
an exponentially large quantum system in its own right, and the questions HSF is
asked (does a sector thermalise? what is the level statistics inside it? what does
an initial state relax to?) all have content. Under $X$ the same rule has
$\sim2^N/\mathrm{poly}$ orbits whose typical length is $O(N)$: exponentially many
``sectors'', each of which is a single trajectory.

\paragraph{It is visible in an entanglement measurement.} Under the Hadamard an
initial computational basis state spreads over its sector, and R8~\S10 measured
exactly this for rule $150$ at $N=11$: the saturation value of the entanglement
entropy clusters by sector, each cluster sitting near the Page value of its own
sector dimension. The fragmentation is therefore an observable statement about a
quantum state. Under $X$ a basis state remains a basis state for all time, so
$S(t)\equiv0$ and there is nothing to measure --- the fragmentation is invisible
to any entanglement probe applied to basis states.

\paragraph{Every \emph{sector} label of the $X$ circuit is classical.} A
permutation of the basis conserves the indicator function of each of its orbits,
so a fragmented automaton has an exponentially large commutant for free. That is
not a discovery about the model; it is the statement that a deterministic map has
trajectories. What R9 finds is the same object in a setting where the sector label
is still a classical function on configurations but the dynamics inside is not ---
and where, by the diagonal-commutant identity of R8~\S10.3, the sector projectors
are precisely the diagonal part of the commutant, i.e.\ everything a classical
partition can see and nothing more. The qualifier matters: the \emph{off}-diagonal
part of the commutant of an X-gate rule with resets is not trivial either, and
R10~\S9 measures it.

\subsection{Why the $X$-gate and the V-free rules are the boring end}

The same three points, run in reverse, and they are worth stating plainly because
the numbers in \S\ref{sec:classes} look impressive on their own.

Under $X$ the $256$ rules split into two kinds and neither supports an interesting
fragmentation claim. For the $16$ reversible rules, sectors \emph{are} orbits,
so the decomposition exists for every bijection whatsoever and contains no
information about locality, constraints or interactions. For the other $240$, the
sectors are basins of attraction: they can be exponentially large, but the
long-time state is a cycle of median length $1$ and a cyclic fraction of
$10^{-7}$, so a sector is an exponentially large funnel into a point. Neither is
Hilbert-space fragmentation in the sense the term is used for constrained quantum
models.

The $81$ V-free rules make the point sharpest, because they are the same circuit
under both gates (Corollary~\ref{cor:nofree}). Whatever \emph{fragmentation} they
show is a property of a classical cellular automaton that a quantum engine
happened to compute. Their \emph{coherence} structure is another matter: rule
$232$ (\rt{DIIE}) is V-free and has a decoherence-free subspace of dimension
$\sim\varphi^N$, confirmed against the exact Cesaro rank
($\dim\operatorname{Fix}(\Phi)=n_{\rm rec}^2$, R10~\S9.4). A classical-population
circuit can protect an exponentially large quantum subspace, and that is not a
classical statement. Four of R9's eight ``fragmentation in an open quantum system'' rules ---
$140$, $196$, $206$, $220$ --- are in this set, so for those four the claim should
be read as classical fragmentation surviving a reset, with the Hadamard playing no
part. The other four ($156$, $198$, $108$, $201$) do use the gate, and for them
the claim stands as stated.

\subsection{Automaton circuits: why the community calls them quantum}
\label{sec:automaton}

The natural objection to R10 is that it studies a classical cellular automaton and
calls it a QCA. The objection is half right, and the half that is wrong is
instructive.

\paragraph{Unitarity is not quantumness.} A permutation matrix is unitary, so a
reversible CA \emph{is} a quantum cellular automaton under the structural
definition --- discrete time, strict locality, a causal cone, a unitary on
$(\mathbb{C}^2)^{\otimes N}$. Nothing is being smuggled. What is true is that
unitarity alone buys none of the phenomena the word ``quantum'' is meant to
signal.

\paragraph{The quantum content is in the initial state, not the gate.} A
permutation circuit maps basis states to basis states, but it is linear, and on a
\emph{superposition} it does not act classically at all. Take $\ket{+}^{\otimes N}$
or any product state off the computational axis: an automaton circuit will
entangle it, and can scramble it as fast as a generic circuit. This is the reason
the class is studied. It sits in the solvable corner --- classically simulable
trajectory by trajectory, so accessible at sizes no exact diagonalisation reaches,
while still supporting entanglement growth, operator spreading and hydrodynamics
that are genuinely quantum questions. Our rule $1$ result in R10~\S5 is of this
type: its recurrent set is exactly the golden-mean shift, which is a statement
about a classical constraint, but it fixes the long-time support of any quantum
state prepared in that sector.

\paragraph{Where the class comes from.} Two lines of work are the relevant
context, and the connection is concrete rather than thematic.

\emph{Rule 54.} The most-studied exactly solvable QCA is the Rule 54 automaton
(studied by, among others, Prosen and collaborators, and by Klobas and Bertini for
its exact steady states and entanglement dynamics). It is a deterministic
reversible cellular automaton whose local update, in the brick-wall formulation,
is $x_i\mapsto x_i\oplus(x_{i-1}\vee x_{i+1})$ --- and that is exactly our
\rt{IVVV}, which our Wolfram numbering calls rule $54$. More generally, for all
$16$ reversible rules the local update is $x_i\mapsto x_i\oplus f(x_{i-1},x_{i+1})$
with $f$ one of the $16$ boolean functions of two arguments, and the induced
elementary CA is the ECA of the \emph{same Wolfram number} --- we checked all $16$.
So the reversible X-gate family is precisely the set of ECAs that are affine in
the centre cell, brick-walled, and Rule 54 is one of its members. R10's structural
results therefore apply to it directly: $X_{54}$ has $2^N/\mathrm{poly}$ orbits of
length $52,55,58,61,64,67$ at $N=19\ldots24$, growing by exactly $3$ per site,
which is the soliton picture on an open chain seen through the sector axis.

\emph{Quantum automaton circuits.} The other line uses exactly this class as a
tool: unitaries that permute the computational basis, deployed to study
fragmentation, freezing and anomalously slow transport at sizes where the
entanglement can still be computed --- work associated with Iadecola and
collaborators, and with the automaton-circuit studies of subdiffusion and
subsystem symmetries. The methodological bargain is explicit there: classical
simulability in exchange for a restricted class of dynamics, with the quantum
questions asked of superposed initial states.

\paragraph{What this report adds to that.} A complete $256$-rule census of the
sector and orbit structure of the class, and one structural statement that the
automaton literature and the constrained-quantum-model literature do not usually
put side by side: the automaton sectors \emph{refine} the Hadamard ones
(Theorem~\ref{thm:refine}). An automaton model is therefore a lower bound on the
fragmentation of its superposing counterpart, never an upper bound, and any
fragmentation found in an automaton circuit should be checked against the
question of whether it survives gluing --- which for $39$ of our $256$ rules it
does not.

\paragraph{A caveat on the attributions.} The two paragraphs above characterise
the literature from memory, without the papers to hand; the identification of
$X_{54}$ with Rule 54 is our own check and is exact, but the names and the
framing attached to the two research lines should be verified against the sources
before this section is used anywhere outside the project.

\section{What each family says that the other cannot}
\label{sec:asym}

The comparison is not symmetric, and the asymmetries are informative.

\paragraph{Only the Hadamard family has branch interference --- which is less than
it sounds.} Destructive cancellation between two branches of a basis state needs
a successor set with more than one element, and under $X$ there is only ever one,
so that specific mechanism is absent. It does not follow that the permutation
family is decoherence-free-trivial, and an earlier draft of both reports said
that it did. A reset has two Kraus operators which agree on populations and
differ on coherences, so the functional graph --- a statement about populations
alone --- cannot see decoherence at all. R10~\S9 computes what it misses: at
$N=12$, $127$ of the $240$ irreversible X rules carry a decoherence-free subspace
of dimension $>1$, and the dimensions are exponential, with rule $1$'s equal to
the golden-mean shift ($F_{N+2}$), rules $76$ and $73$'s to the no-\rt{111} shift
(tribonacci), and rule $22$'s to the Lucas numbers $L_{N+1}$. The X-gate family
is therefore a classical control \emph{for populations only}. That is still the
useful reading --- \S\ref{sec:classes} shows the sector axis is almost entirely
graph combinatorics --- but the coherence axis is not classical at all.

\paragraph{Only the permutation family has a clean recurrent structure.} With
out-degree one, $n_\rec=n_\wcc$ identically and every basin is a sector, so
questions about attractors have exact answers (R10~\S\S5--8). In the Hadamard
family the same questions are genuinely harder --- rule $22$ has three attractors
inside two sectors --- which is why R9 could assert A4 in one direction only. The
permutation family is where the attractor-level statements can be proved.

\paragraph{The exclusion curve belongs to both.} $ab\ge2$ is a theorem about
partitions and applies wherever sectors partition the basis, so it binds both
families identically, and neither produced a violation. It does \emph{not} apply
to cycles, and R10~\S5 makes that precise: the recurrent set is not the whole
space, so the analogous constraint degenerates.

\section{Open items}

\begin{enumerate}
\item \textbf{\rt{pbc}.} Held for both families. The refinement theorem is proved
  without reference to the boundary, so it should carry over verbatim, and it will
  be worth checking that the ring does not break the $81$-rule identity.
\item \textbf{The gluing map.} Theorem~\ref{thm:refine} says each Hadamard sector
  is a union of $X$ sectors, but nothing here counts them: the natural next
  observable is the number of $X$ sectors per Hadamard sector, whose growth would
  measure exactly how much the non-flip branch glues.
\item \textbf{A third gate.} $X$ and $H$ are the two extremes --- a permutation and
  a maximally superposing gate. A partial rotation would interpolate, and the
  refinement theorem suggests the sector lattice would interpolate monotonically
  with it.
\item \textbf{The four classical rules in R9's open-fragmented list.} $140$,
  $196$, $206$, $220$ deserve a purely combinatorial derivation of their $2584$
  sectors; nothing quantum is involved and the count is Fibonacci.
\end{enumerate}

\section{Reproduce}

\begin{verbatim}
cd code
python -m qca_fragmentation.permutation.compare --bc obc0
python -m qca_fragmentation.permutation.compare_tables --bc obc0
python -m qca_fragmentation.permutation.compare_figures --bc obc0 --rebuild
python -m pytest qca_fragmentation/tests/test_gate_compare.py -q
\end{verbatim}

\begin{sloppypar}
Code: \fn{permutation/compare.py} (the refinement checks, the common-window
descriptors, the correlations), \fn{permutation/compare\_tables.py},
\fn{permutation/compare\_figures.py}. Data: \fn{results\_wcc/} and
\fn{results\_xgate/}, \fn{analytics/gate\_compare\_obc0.json}.
\end{sloppypar}

\end{document}
